% GERADO por tools/generation/gerar_secoes_latex.py a partir de
% artifacts/paper/secoes/08-discussao-limitacoes.md. Nao editar a mao: a proxima geracao
% desfaz. Corrija o Markdown e regere.
\textbf{A subdeterminação não se manifesta como contradição, e é por isso que ela escapa.} Nenhum dos desvios contraria uma frase do modelo de referência; cada um decorre de decisão que o texto não toma e que quem formaliza tem de tomar. O ponto de abertura tem forma reconhecível: um verbo de relação que admite mais de uma leitura ontológica, uma enumeração de coisas com identidades distintas tratada como categoria única, um mesmo fenômeno nomeado como componente e como ação sem que se diga qual dos dois vira elemento formal. Reconhecer essa forma é o que o procedimento oferece a quem for auditar outro modelo em prosa: não uma lista de defeitos a procurar, mas os lugares onde eles nascem.

\textbf{As ferramentas atestam conformidade, e conformidade não é adequação.} O verificador de conformidade à UFO nada encontrou na formalização defeituosa, e nada encontraria: as deficiências mais consequentes são defeitos de significado, fora do alcance de qualquer regra sintática. A lição é de método: um relatório vazio só é interpretável se alguém provar que o instrumento dispara, e o limite de cada um precisa ser declarado junto com o resultado. A recíproca vale, e um relator ternário que duas leituras manuais deixaram passar foi achado pela ferramenta: a análise humana erra por familiaridade com o texto, a ferramenta por não o ler, e é a combinação que sustenta o resultado.

\textbf{A instanciação expõe o que o argumento não expõe.} As questões de competência foram feitas para validar, e o que fizeram de mais útil foi diagnosticar: revelaram que não há vínculo entre a observação e o projeto que ela concerne, que a proveniência para no gerenciamento responsável e que a motivação só é dizível como anotação --- nenhuma delas apareceu na leitura do modelo nem nos relatórios das ferramentas.

\textbf{O que isso diz sobre modelos de referência descritos em prosa.} As definições existiam: o apêndice da conceituação define 53 termos, e nenhum chegou ao modelo formal. Publicá-lo sem os artefatos intermediários que o produziram devolve ao leitor a tarefa de reconstruir as decisões, e é aí que as interpretações divergem. Publicar a conceituação ao lado do modelo fecha boa parte da abertura; dizer o tipo ontológico de cada relação enunciada fecha o resto.

Seis limitações precisam ser declaradas. \textbf{Não houve validação com especialistas}, e a ausência pesa contra um pano de fundo específico: o modelo analisado foi avaliado, na literatura, por grupos focais e por estudos de caso múltiplos publicados nos Anais do SBSI \cite{vieira2022evaluating,vieira2023survey}. Esse é o padrão do domínio, e este trabalho não o cumpre: o que oferece no lugar atesta capacidade de resposta e conformidade, não aceitação pela comunidade. \textbf{Há um formalizador}, e superá-lo exige formalizações independentes (\S{}9). \textbf{O segundo observatório da comparação é sintético}, e a assimetria é em parte por construção.

\textbf{O artefato tem lacunas registradas}, nenhuma corrigida em silêncio: a camada de atributos que a conceituação especifica não existe no modelo; um relator segue candidato a evento; e o coletivo do observatório reúne membros de naturezas incompatíveis. A \textbf{verificação de consistência é do perfil OWL 2 RL}, e ele não avalia cardinalidade qualificada em posição de superclasse --- o que a transformação gera das multiplicidades. Por fim, \textbf{ficaram fora de escopo, por decisão de pesquisa}, a implantação operacional em produção, a axiomatização pesada com regras e cadeias de papéis, e o alinhamento a ontologias externas.
